% ============================================================
%  ch2_I_b  第一章末（原书 pp. 10--11，即 PDF 第 11--12 页）
% ============================================================

\origpage{10}

\par\medskip
\noindent\textbf{所需的材料：}

\begin{itemize}
\item Arthur--Selberg 迹公式；
\item chtoucas 的阿代尔描述（Drinfeld）；
\item 基本引理（Drinfeld 的球函数）。
\end{itemize}

\par\medskip
\noindent\textbf{非常重要的注：}\ 主项不依赖于 $p$。它恰好对应所寻求的表示

\[ \bigoplus_{\pi} \pi \otimes q'^*\sigma_\pi \otimes q''^*\wt{\sigma}_\pi(1 - r). \]

\par\medskip
\noindent\textbf{随之而来的新问题：}\ 开子叠 $\Cht^{r,\ol{p}\le p}_N/a^{\ZZ}$ 不被 Hecke 对应 $f \in \cH^r_N$ 稳定。

不动点公式没有上同调意义（单位对应 $\mathbb{I}_N$ 除外）。

\medskip
\noindent\textbf{想法：}\ 构造 $\Cht^{r,\ol{p}\le p}_N/a^{\ZZ}$ 的紧化 $\ol{\Cht}^{r,\ol{p}\le p}_N/a^{\ZZ}$，使得可以用归一化延拓对应 $f$，并在上同调上重新得到一个作用。

$\rightarrow$ 当 $N = \emptyset$ 时，将构造 $X \times X$ 上紧合且光滑的 $\ol{\Cht}^{r,\ol{p}\le p}/a^{\ZZ}$。\\
当 $N \neq \emptyset$ 时，情况会复杂得多 $\ldots$

\medskip
\noindent\textbf{新问题：}\ 经归一化延拓后的 Hecke 对应各自作用在上同调上，但它们并不确定出一个 $\cH^r_N$ 的作用。

\origpage{11}

\par\medskip
\noindent\textbf{总结起来：}

\begin{center}
\renewcommand{\arraystretch}{1.35}
\begin{tabular}{p{4.4cm}p{4.4cm}p{4.4cm}}
\toprule
\multicolumn{1}{c}{chtoucas} & \multicolumn{1}{c}{截断} & \multicolumn{1}{c}{紧化}\\
\midrule
$\Cht^r_N/a^{\ZZ}$ \ $\xleftarrow{\quad\quad}$ & $\Cht^{r,\ol{p}\le p}_N/a^{\ZZ}$ \ $\xrightarrow{\quad\quad}$ & $\ol{\Cht}^{r,\ol{p}\le p}_N/a^{\ZZ}$\\[2pt]
\multicolumn{1}{c}{$\vdots$} & \multicolumn{1}{c}{$\vdots$} & \multicolumn{1}{c}{$\vdots$}\\[2pt]
$\bullet$ 维数 $2r-2$ 的光滑态射到 $(X - N) \times (X - N)$
& $\bullet$ 有限型
& $\bullet$ 在 $X \times X$ 上紧合且光滑（若 $N = \emptyset$）\\
$\bullet$ $\cH^r_N$ 的作用
& $\bullet$ 不动点计数
& $\bullet$ 对应作用在上同调上\\
\midrule
\multicolumn{1}{c}{\textbf{但是}} & \multicolumn{1}{c}{\textbf{但是}} & \multicolumn{1}{c}{\textbf{但是}}\\
\midrule
$\bullet$ 非有限型
& $\bullet$ 不再有 Hecke 对应
& $\bullet$ 没有 $\cH^r_N$ 的作用；\newline $\bullet$ 没有不动点计数\\
\bottomrule
\end{tabular}
\end{center}

\par\medskip
\noindent\textbf{想法：}\ 在这三个对象的 $\ell$-进上同调中分离出一块，它在三者中是同一块，并把它算出来。

\begin{definition*}
\emph{$G_{F^2}$ 的不可约 $\ell$-进表示称为 \textbf{$r$-可忽略的}，若它是形如}
\[ q'^*\sigma_1 \otimes q''^*\sigma_2 \qquad (q', q'' : X \times X \rightrightarrows X) \]
\emph{的表示的直和项，其中 $\sigma_1, \sigma_2$ 是 $G_F$ 的两个维数 $< r$ 的不可约表示。否则就称为\textbf{本质的}。}
\end{definition*}

\par\medskip
\noindent\textbf{我们将证明：}\ 这三个对象有相同的本质上同调。它是

\[ \bigoplus_{\pi \in \{\pi\}^N_r} \pi \otimes q'^*\sigma_\pi \otimes q''^*\wt{\sigma}_\pi(1 - r). \]
